% EXPECTED-LEAKS: affiliation, email, institution, city, country, orcid, settopmatter, printacmref, copyrightyear, acmConference, ccsdesc
\documentclass[sigconf]{acmart}

\settopmatter{printacmref=false}
\copyrightyear{2026}
\acmConference[SIGMOD '26]{ACM SIGMOD 2026}{June 2026}{Seattle, WA}

\begin{document}
\title{Learned Cardinality Estimation at Scale}

\author{Carol Davis}
\affiliation{%
  \institution{Example Corp}
  \city{San Francisco}
  \country{USA}
}
\email{carol@example.com}

\begin{abstract}
Cardinality estimation is a long-standing database problem. We propose a learned model that reduces estimation error by 40\% on industrial workloads.
\end{abstract}

\begin{CCSXML}
<ccs2012>
<concept><concept_id>10002951.10002952.10003219</concept_id>
<concept_desc>Information systems~Query optimization</concept_desc></concept>
</ccs2012>
\end{CCSXML}

\ccsdesc[500]{Information systems~Query optimization}

\keywords{cardinality, learned indexes, query optimization}

\maketitle

\section{Introduction}
Modern databases rely on accurate cardinality estimates to choose efficient query plans. Traditional histogram-based approaches break down under skew and correlated predicates.

\section{Model Architecture}
Our neural estimator uses an encoder-decoder topology with attention over predicate columns.

\section{Experimental Results}
We evaluate on IMDB, TPC-H, and a proprietary dataset. Table results show 40\% reduction in median q-error.

\section{Conclusion}
Learned models complement traditional statistics and warrant deeper production study.

\end{document}
